\section{官方演练观察：结果与可识别的瓶颈}
\label{field:sec-official}
用户提供了2026年9月12日在官方模拟器完成的Q3、Q4各10轮演练：每局包含脱敏请求日志、控制器摘要及模拟器退出后结果文件，案例关联由\texttt{manifest.csv}固定。执行时采用原\texttt{refined}，不是新增\texttt{field}。事后公开的源数和类型仅用于评价，绝不传入在线规划器。\texttt{official\_analysis.py}逐请求与响应独立复算移动距离、检测、换频及清除费用，并按微秒累计与每步虚拟时刻核对；客户端派生的距离和换频字段只用于交叉检查。

\begin{table}[H]\centering\small
\caption{20轮官方演练的汇总；每题10轮均全清，时间单位秒/源}
\begin{tabular}{lrrrrrr}
\toprule
问题 & 局数 & 全清源数 & 局等权T/K & 汇总T/K & 移动占比/\% & 检测/局 \\
\midrule
Q3 & 10 & 139 & 238.00 & 235.82 & 74.38 & 128.90 \\
Q4 & 10 & 138 & 531.18 & 519.38 & 73.20 & 311.80 \\
\bottomrule
\end{tabular}

\end{table}
其中“局等权”为$\frac1{10}\sum_j T_j/K_j$，“汇总”为$\sum_jT_j/\sum_jK_j$，后者给予源数更多的局更高权重。问题三139个、问题四138个源均清除；局等权均值分别为\OfficialQThreeMeanRun{}、\OfficialQFourMeanRun{}秒/源。问题四第8局只有10个源，仍需完成25点搜索，因此$T/K=768.87$秒/源，说明小源数案例的固定发现成本不能忽略。

\begin{table}[H]\centering\small
\caption{20轮官方演练逐局结果；$N$为退出后标签，不是正式测试分母}
\begin{tabular}{rrlrrrr}
\toprule
题号 & 轮次 & 案例编码 & $N$ & $K$ & $T$/s & $T/K$/(s/源) \\
\midrule
Q3 & 1 & BUWJ-NM9Q-E3FP-KBNQ & 16 & 16 & 3501.15 & 218.82 \\
Q3 & 2 & XSTE-PSFN-9F8U-8P2D & 12 & 12 & 2929.34 & 244.11 \\
Q3 & 3 & WQXS-VYHK-4B2N-4NWU & 13 & 13 & 3230.83 & 248.53 \\
Q3 & 4 & 3YB4-CC66-WNAD-XVQ9 & 15 & 15 & 3301.25 & 220.08 \\
Q3 & 5 & AKGC-XRSQ-KQ42-DE4U & 14 & 14 & 2867.24 & 204.80 \\
Q3 & 6 & F22D-5NT7-9XRW-7DY6 & 13 & 13 & 3119.85 & 239.99 \\
Q3 & 7 & FDS3-ZJ7Q-V3E9-2QDN & 11 & 11 & 3334.99 & 303.18 \\
Q3 & 8 & XMZP-GEJC-W9KA-UTWY & 14 & 14 & 3452.13 & 246.58 \\
Q3 & 9 & 8V3W-N88T-ED46-PXCE & 15 & 15 & 3294.33 & 219.62 \\
Q3 & 10 & HGVD-7ZTC-XNGP-2JRT & 16 & 16 & 3748.24 & 234.26 \\
Q4 & 1 & ZGME-Z9ZY-UQDJ-UZRF & 15 & 15 & 7091.90 & 472.79 \\
Q4 & 2 & AQVF-W7A2-YAFY-F6JT & 16 & 16 & 7385.70 & 461.61 \\
Q4 & 3 & HS4Z-3F9D-BE4W-FUBW & 13 & 13 & 7233.66 & 556.44 \\
Q4 & 4 & BFSU-HJRZ-2WQX-63MP & 14 & 14 & 7461.23 & 532.95 \\
Q4 & 5 & RFFB-U23B-VMS7-ASRE & 15 & 15 & 6420.23 & 428.02 \\
Q4 & 6 & AJFX-AY9F-VETM-UBJ8 & 15 & 15 & 7089.98 & 472.67 \\
Q4 & 7 & CR6D-5DV5-WCJQ-6N3F & 16 & 16 & 7677.93 & 479.87 \\
Q4 & 8 & TY94-4B9M-J2RU-EDHA & 10 & 10 & 7688.73 & 768.87 \\
Q4 & 9 & 5J2D-BATC-W2P2-BNQ2 & 13 & 13 & 7151.24 & 550.10 \\
Q4 & 10 & CENN-WQ49-E4D5-9JJ7 & 11 & 11 & 6473.24 & 588.48 \\
\bottomrule
\end{tabular}

\end{table}

移动约占Q3的74.38\%、Q4的73.20\%，失败光学清除仅占0.11\%、0.12\%。因此优先改进移动与观测价值，而不是以减少失败尝试次数这一小分项代替总时间目标。20轮未出现\texttt{near}、HTTP拒绝或传输重试；它们支持正常协议联通，不支持“官方已经验证故障恢复”的结论。Q3、Q4在已发现频道后分别仍有223和530次无信号反馈，尤其不能对Q4据此删去1000米位置圆盘。

\paragraph{严格回放与反事实边界。}
历史日志没有真实源位置、接收半径或发射方向；成功清除点距源至多20米，不能作为源坐标真值。严格回放只在新请求与记录路径、频道、坐标一致时返回原反馈，一旦分岔立即停止，不为未测位置猜反馈。原算法在本机可完整匹配8/20局，其余在任务顺序或对称测点选择处分岔；这与跨平台浮点近似及近等价几何决策相容，但仅凭日志不能唯一确定原因。20局的实际费用仍全部由原始坐标和响应核验，不依赖策略回放成功。故以下新策略收益来自独立合成配对，不是这20局的官方反事实成绩；新\texttt{field}仍需另做官方演练。
